% Pursuing Stacks (cotes 134-2 to 134-8): archive cote 134-7, site page 33 (PDF page 34), typescript page 475.
% typescript. Draft compared: ps6.tex:1543-1620.
% AI-transcribed by Claude Opus 5.5 (claude-opus-5-5, adaptive thinking, effort high), Message Batches API, batch ps3 (msgbatch_01XdFmuVVBAcgc7EFxdtCqsj), 2026-11-12; reviewed against both supplied scan halves; see CHECK_DECISIONS.tsv.
\scanpage{134-7/33}
\mspage{475}

We'll more specifically work in the category $U(k)^\hand{\bullet}$ of \uline{pointed} objec\unsure{ts}
of $U(k)$, namely objects endowed with a section over the final object $e$,
the so-called \uline{pointed} \uline{unipotent bundles}. We now have a functor
\[ W^\hand{\bullet} : \quad Ab_k \hand{\longrightarrow} U(k)^\hand{\bullet} \tag{L} \]
deduced from $W$ using the fact that $W(0)=e$ , and a ``\uline{pointed linearizati-}%
\uline{on functor}''
\[ L^\hand{\bullet} \text{or } L_{pt} \;:\; U(k)^\hand{\bullet} \hand{\longrightarrow} Ab_k \quad , \tag{M} \]
giving rise to a commutative diagram similar to (J')
% Hand-drawn arrows are marked blue; typed objects remain plain.
\diagram{
\begin{tikzcd}
Ab_k \arrow[color=blue,rr, "W^\hand{\bullet}"] \arrow[color=blue,dr, "\hand{\Gamma}_k"'] & & U(k)^\hand{\bullet} \arrow[color=blue,dl, "L^\hand{\bullet}"] \\
 & Ab_k & ;
\end{tikzcd}
\tag{M'}
}
the notation $L^\hand{\bullet}$ seems here the most coherent one, but may bring about
confusion with the similar notation for some cochain complex say, there-%
fore we had first used the alternative notation $L_{pt}$ , to which one may
still come back if needed. This time the functor $L^\hand{\bullet}$ commutes to ring
extension without any grain of salt. We have of course a canonical
embedding
\[ L^\hand{\bullet}(X) \hand{\hookrightarrow} L(X) \tag{N} \]
\deleted{corresponding to} \inserted{defined via} the corresponding embedding for an object $M$ of $Ab_k$
\[ \hand{\Gamma}_k(M) \hand{\hookrightarrow} \hand{\Gamma}_k{}^{\hand{\wedge}}(M) \quad , \tag{N'} \]
by which $L(X)$ may be viewed as the completion of $L^\hand{\bullet}(X)$ with respect to
the topology it induces on it, which is a \uline{canonical} \uline{topology} on \deleted{\unsure{X|it}}\hand{$L^\bullet(X)$}.
% NOTE: a typed ``X'' and ``Y'' stand between this line and the next, under ``the'' and ``topology''
$X$ \quad $Y$
Maps in $U(k)^\hand{\bullet}$ correspond to \uline{k-linear} \uline{maps}
\[ L^\hand{\bullet}(X) \hand{\longrightarrow} L^\hand{\bullet}(Y) \]
which are moreover \uline{continous}, \inserted[typed]{and} \uline{compatible with} \inserted[typed]{co}\uline{augmentation}\deleted{\unsure{s}}\xout{as well as}
(i.e. transforms 1 into 1) \uline{as well as with augmentations and diagonal}
\uline{maps}. I wonder if there is any simple \hand{c}haracterization of submersions
\inserted[typed]{in $U(k)^\hand{\bullet}$} in terms of the corresponding map between the linearizations. At any
rate, an object \inserted[typed]{$X$} of $U(k)^\hand{\bullet}$ is smooth resp. flat iff $L^\hand{\bullet}(X)$ is a projecti-%
ve resp. a flat $k$-module.

For any natural integer, we want now to define a model category
\[ \hand{$\mathcal{M}$}_n(k) \hand{$\subset$} \uline{Hom}(\hand{$\Delta^{\circ}$}, U(k)^\hand{\bullet}) \quad , \tag{O} \]
which should be a full subcategory of the category of semisimplicial
\deleted{\xout{objects in U(k)}}

\margin{33}

% Visual collation: both supplied image halves inspected; see CHECK_DECISIONS_134_7_8.tsv.
